% Options for packages loaded elsewhere
\PassOptionsToPackage{unicode}{hyperref}
\PassOptionsToPackage{hyphens}{url}
\PassOptionsToPackage{dvipsnames,svgnames,x11names}{xcolor}
%
\documentclass[
  10pt,
]{article}
\usepackage{amsmath,amssymb}
\usepackage{iftex}
\ifPDFTeX
  \usepackage[T1]{fontenc}
  \usepackage[utf8]{inputenc}
  \usepackage{textcomp} % provide euro and other symbols
\else % if luatex or xetex
  \usepackage{unicode-math} % this also loads fontspec
  \defaultfontfeatures{Scale=MatchLowercase}
  \defaultfontfeatures[\rmfamily]{Ligatures=TeX,Scale=1}
\fi
\usepackage{lmodern}
\ifPDFTeX\else
  % xetex/luatex font selection
\fi
% Use upquote if available, for straight quotes in verbatim environments
\IfFileExists{upquote.sty}{\usepackage{upquote}}{}
\IfFileExists{microtype.sty}{% use microtype if available
  \usepackage[]{microtype}
  \UseMicrotypeSet[protrusion]{basicmath} % disable protrusion for tt fonts
}{}
\makeatletter
\@ifundefined{KOMAClassName}{% if non-KOMA class
  \IfFileExists{parskip.sty}{%
    \usepackage{parskip}
  }{% else
    \setlength{\parindent}{0pt}
    \setlength{\parskip}{6pt plus 2pt minus 1pt}}
}{% if KOMA class
  \KOMAoptions{parskip=half}}
\makeatother
\usepackage{xcolor}
\usepackage[margin=2.3cm]{geometry}
\usepackage{longtable,booktabs,array}
\usepackage{calc} % for calculating minipage widths
% Correct order of tables after \paragraph or \subparagraph
\usepackage{etoolbox}
\makeatletter
\patchcmd\longtable{\par}{\if@noskipsec\mbox{}\fi\par}{}{}
\makeatother
% Allow footnotes in longtable head/foot
\IfFileExists{footnotehyper.sty}{\usepackage{footnotehyper}}{\usepackage{footnote}}
\makesavenoteenv{longtable}
\setlength{\emergencystretch}{3em} % prevent overfull lines
\providecommand{\tightlist}{%
  \setlength{\itemsep}{0pt}\setlength{\parskip}{0pt}}
\setcounter{secnumdepth}{-\maxdimen} % remove section numbering
\newlength{\cslhangindent}
\setlength{\cslhangindent}{1.5em}
\newlength{\csllabelwidth}
\setlength{\csllabelwidth}{3em}
\newlength{\cslentryspacingunit} % times entry-spacing
\setlength{\cslentryspacingunit}{\parskip}
\newenvironment{CSLReferences}[2] % #1 hanging-ident, #2 entry spacing
 {% don't indent paragraphs
  \setlength{\parindent}{0pt}
  % turn on hanging indent if param 1 is 1
  \ifodd #1
  \let\oldpar\par
  \def\par{\hangindent=\cslhangindent\oldpar}
  \fi
  % set entry spacing
  \setlength{\parskip}{#2\cslentryspacingunit}
 }%
 {}
\usepackage{calc}
\newcommand{\CSLBlock}[1]{#1\hfill\break}
\newcommand{\CSLLeftMargin}[1]{\parbox[t]{\csllabelwidth}{#1}}
\newcommand{\CSLRightInline}[1]{\parbox[t]{\linewidth - \csllabelwidth}{#1}\break}
\newcommand{\CSLIndent}[1]{\hspace{\cslhangindent}#1}
\ifLuaTeX
\usepackage[bidi=basic]{babel}
\else
\usepackage[bidi=default]{babel}
\fi
\babelprovide[main,import]{spanish}
% get rid of language-specific shorthands (see #6817):
\let\LanguageShortHands\languageshorthands
\def\languageshorthands#1{}
\ifLuaTeX
  \usepackage{selnolig}  % disable illegal ligatures
\fi
\IfFileExists{bookmark.sty}{\usepackage{bookmark}}{\usepackage{hyperref}}
\IfFileExists{xurl.sty}{\usepackage{xurl}}{} % add URL line breaks if available
\urlstyle{same}
\hypersetup{
  pdftitle={Publicaciones previstas --- CorAnTool},
  pdfauthor={Santiago Casanova-Arenillas},
  pdflang={es},
  colorlinks=true,
  linkcolor={Maroon},
  filecolor={Maroon},
  citecolor={Blue},
  urlcolor={Blue},
  pdfcreator={LaTeX via pandoc}}

\title{Publicaciones previstas --- CorAnTool}
\usepackage{etoolbox}
\makeatletter
\providecommand{\subtitle}[1]{% add subtitle to \maketitle
  \apptocmd{\@title}{\par {\large #1 \par}}{}{}
}
\makeatother
\subtitle{Cartera, métodos, defensa ante revisores y borradores}
\author{Santiago Casanova-Arenillas}
\date{01-11-2026 · software V29}

\begin{document}
\maketitle

{
\hypersetup{linkcolor=}
\setcounter{tocdepth}{2}
\tableofcontents
}
\hypertarget{publicaciones-previstas-corantool-antes-ichantool}{%
\section{Publicaciones previstas --- CorAnTool (antes
IchAnTool)}\label{publicaciones-previstas-corantool-antes-ichantool}}

Borradores, revistas candidatas y estado de cada trabajo. Estado al
\textbf{01-11-2026} (software en la entrega V28). Todo lo marcado
\texttt{{[}TODO{]}} está pendiente; ningún resultado de estos borradores
es definitivo.

\begin{quote}
Aviso: si el repositorio es \textbf{público}, no subir aquí datos ni
resultados inéditos de la tesis hasta que el manuscrito esté enviado o
publicado como preprint (o pasar los borradores a un repositorio
privado).
\end{quote}

\hypertarget{cartera-de-publicaciones}{%
\subsection{Cartera de publicaciones}\label{cartera-de-publicaciones}}

\begin{longtable}[]{@{}
  >{\raggedright\arraybackslash}p{(\columnwidth - 8\tabcolsep) * \real{0.0222}}
  >{\raggedright\arraybackslash}p{(\columnwidth - 8\tabcolsep) * \real{0.3333}}
  >{\raggedright\arraybackslash}p{(\columnwidth - 8\tabcolsep) * \real{0.3333}}
  >{\raggedright\arraybackslash}p{(\columnwidth - 8\tabcolsep) * \real{0.1778}}
  >{\raggedright\arraybackslash}p{(\columnwidth - 8\tabcolsep) * \real{0.1333}}@{}}
\toprule\noalign{}
\begin{minipage}[b]{\linewidth}\raggedright
\#
\end{minipage} & \begin{minipage}[b]{\linewidth}\raggedright
Trabajo y aporte
\end{minipage} & \begin{minipage}[b]{\linewidth}\raggedright
Revista principal (alternativas)
\end{minipage} & \begin{minipage}[b]{\linewidth}\raggedright
Depende de
\end{minipage} & \begin{minipage}[b]{\linewidth}\raggedright
Estado
\end{minipage} \\
\midrule\noalign{}
\endhead
\bottomrule\noalign{}
\endlastfoot
1 & \textbf{Software: CorAnTool.} Herramienta abierta y reproducible &
JOSS (Applied Computing and Geosciences; Computers \& Geosciences) &
Release con DOI & Borrador en inglés \\
2 & \textbf{Dataset anotado y segmentación automática.} Datos abiertos y
método & Earth System Science Data (Scientific Data; AI in Geosciences;
Computers \& Geosciences) & 1 e imágenes con permiso & Esquema \\
3 & \textbf{Validación frente a índices visuales y aplicación
paleoambiental.} Interpretación geológica & Palaeogeography,
Palaeoclimatology, Palaeoecology (Ichnos; Marine Geology; Sedimentary
Geology) & 1 (y 2) & Esquema \\
4 & \textbf{Clasificación automática de intervalos con datos
multiproxy.} Método de aprendizaje automático & Computers \& Geosciences
(AI in Geosciences; Applied Computing and Geosciences; Marine Geology) &
1 y sondeos con eventos descritos & Esquema y primeras pruebas \\
5 & \textbf{Nota breve en español.} Difusión en España & Geogaceta
(Spanish Journal of Palaeontology) & 1 & Idea \\
\end{longtable}

\textbf{Orden recomendado:} 1 $\rightarrow$ 5 $\rightarrow$ 4 $\rightarrow$ 2 $\rightarrow$ 3.

\begin{itemize}
\tightlist
\item
  El 1 da el DOI que citan los demás; el 5 se puede escribir a la vez
  con el mismo material.
\item
  El 4 es el más avanzado en software (V26--V27: familias,
  clasificadores, red neuronal, validación) y solo necesita datos con
  eventos descritos de forma independiente.
\item
  El 2 y el 3 dependen de trabajo con otras personas (anotadores,
  observadores a ciegas) y llevan más tiempo.
\end{itemize}

\hypertarget{quuxe9-puede-afirmar-hoy-cada-uno-resumen}{%
\subsection{Qué puede afirmar hoy cada uno
(resumen)}\label{quuxe9-puede-afirmar-hoy-cada-uno-resumen}}

\begin{longtable}[]{@{}
  >{\raggedright\arraybackslash}p{(\columnwidth - 4\tabcolsep) * \real{0.0270}}
  >{\raggedright\arraybackslash}p{(\columnwidth - 4\tabcolsep) * \real{0.4865}}
  >{\raggedright\arraybackslash}p{(\columnwidth - 4\tabcolsep) * \real{0.4865}}@{}}
\toprule\noalign{}
\begin{minipage}[b]{\linewidth}\raggedright
\#
\end{minipage} & \begin{minipage}[b]{\linewidth}\raggedright
Se puede afirmar
\end{minipage} & \begin{minipage}[b]{\linewidth}\raggedright
No se puede afirmar todavía
\end{minipage} \\
\midrule\noalign{}
\endhead
\bottomrule\noalign{}
\endlastfoot
1 & Reproduce exactamente los resultados de 2020; es determinista;
integra imágenes, registros, intervalos, clasificación, mapas y figuras
& Rendimiento de la detección automática en datos reales \\
2 & Las herramientas de anotación existen y producen máscaras abiertas &
Precisión de la segmentación (faltan datos anotados) \\
3 & La app registra BI/ii manuales junto a los automáticos &
Concordancia con observadores (falta el experimento) \\
4 & La cadena funciona con verdad conocida (sintético: RF $\approx$ 91 \%, red $\approx$
80 \%) & Capacidad predictiva real: con los datos de la tesis, 18 \% de
acierto equilibrado dejando un sondeo fuera (7 clases) \\
5 & Funcionalidad y ejemplo con datos públicos & --- \\
\end{longtable}

\hypertarget{documentaciuxf3n-de-muxe9todos}{%
\subsection{Documentación de
métodos}\label{documentaciuxf3n-de-muxe9todos}}

Qué algoritmos usa cada artículo, qué se puede afirmar, qué falta
validar y respuestas preparadas a revisores: METODOS\_Y\_DEFENSA.md.
Detalle técnico con fórmulas: ../docs/DOCUMENTACION\_TECNICA.md.
Revistas: journals.md.

\hypertarget{antes-de-enviar}{%
\subsection{Antes de enviar}\label{antes-de-enviar}}

\begin{itemize}
\tightlist
\item
  Confirmar coautorías (F.J. Rodríguez-Tovar, F. Martínez-Ruiz, J.
  Dorador u otros según aporten datos o validación).
\item
  ORCID de todos los autores.
\item
  Release con DOI de Zenodo del código citado (indicar la entrega,
  p.~ej. V28).
\item
  Comprobar en la web de cada revista: alcance, formato, APC y política
  de preprints.
\item
  Preprint opcional en EarthArXiv.
\end{itemize}

\newpage

\hypertarget{muxe9todos-y-defensa-ante-revisores}{%
\section{Métodos y defensa ante
revisores}\label{muxe9todos-y-defensa-ante-revisores}}

Estado al 01-11-2026 (software en la entrega \textbf{V28}, aplicación
web 2.0.1). Cartera completa en README.md.

\hypertarget{lo-que-hay-que-saber-del-software-antes-de-escribir-cualquier-artuxedculo}{%
\subsection{Lo que hay que saber del software antes de escribir
cualquier
artículo}\label{lo-que-hay-que-saber-del-software-antes-de-escribir-cualquier-artuxedculo}}

\begin{itemize}
\tightlist
\item
  \textbf{Nombre:} desde V16 el programa se llama \textbf{CorAnTool}
  (antes IchAnTool). En los artículos: «CorAnTool (formerly IchAnTool;
  Casanova-Arenillas et al., 2020)». Citar la entrega usada (p.~ej.
  «V28») y, al enviar, la release con DOI.
\item
  \textbf{Icnotaxones (imágenes):} tres modelos (sondeo, proyecto, todos
  los proyectos; §10.6 bis) y tres métodos (k-NN, bosque aleatorio,
  Bayes ingenuo; §10.6 ter). Desde V22 la evaluación agrupa por imagen
  (sección) y cuenta como error cualquier predicción de un taxón ausente
  del conjunto de prueba. Hay que decir qué modelo y método se usó y su
  acierto equilibrado. Las «probabilidades» de RF, NB y k-NN son
  proporciones de voto o verosimilitudes normalizadas, no probabilidades
  calibradas.
\item
  \textbf{Multiproxy (V26--V27):} familias de intervalos (p.~ej. periodo
  frío/cálido; baja oxigenación con sapropel/ORL) que se clasifican a
  mano en cualquier sondeo, y \textbf{clasificadores definidos por el
  usuario} (qué familias, con qué variables y familias de entrada, qué
  modelo: bosque aleatorio, red neuronal MLP, k-NN, Bayes, k-medias),
  validados dejando un sondeo fuera (§21.4--§21.5). Solo se puntúan
  intervalos descritos, no los etiquetados por reglas.
\item
  \textbf{Selección de modelo:} comparar varios clasificadores con la
  misma validación y quedarse con el mejor \textbf{sesga el acierto al
  alza}. Un artículo debe informar de todos los modelos probados o
  reservar sondeos para una prueba final independiente (pendiente de
  implementar en la interfaz).
\item
  \textbf{Sondeos compartidos entre proyectos (V23):} cuentan una sola
  vez en mapas y análisis entre proyectos, y los ejemplos de icnotaxones
  repetidos se descartan. La coincidencia es por nombre: revisar los
  marcados «en N proyectos».
\item
  \textbf{Mapas:} los valores usan solo datos reales; la interpolación
  (IDW, kriging, TIN, RBF\ldots; §21.6 ter) mide distancias en km y es
  exploratoria. Los mapas de z por sondeo eliminan el gradiente absoluto
  entre sondeos: úsense para patrones relativos. Los fondos (Esri,
  GEBCO) son contexto visual y se citan en el pie.
\item
  \textbf{Datos de la tesis:} LC06, LC14 y LC20 tienen profundidad
  supuesta (secciones apiladas, §21.8 ter): no usarlos en una figura
  publicada sin verificarla. Las cifras obtenidas con versiones
  anteriores a V22 deben recalcularse.
\end{itemize}

Para cada artículo previsto: qué partes de IchAnTool usa, qué se puede
afirmar con lo que ya está hecho, qué falta comprobar antes de enviarlo
y las preguntas que probablemente hará un revisor, con una respuesta
preparada.

Los detalles de cada algoritmo (fórmulas, parámetros, archivos de
código) están en \texttt{docs/DOCUMENTACION\_TECNICA.md}. Aquí se cita
la sección (§). El historial de cambios está en
\texttt{docs/HISTORIAL\_VERSIONES.md}.

Regla general: \textbf{solo se afirma lo que está medido}. Lo que no
está validado se presenta como limitación o trabajo futuro.

\begin{center}\rule{0.5\linewidth}{0.5pt}\end{center}

\hypertarget{artuxedculo-1-software-corantool-joss}{%
\subsection{Artículo 1 --- Software: CorAnTool
(JOSS)}\label{artuxedculo-1-software-corantool-joss}}

\hypertarget{quuxe9-describe}{%
\subsubsection{Qué describe}\label{quuxe9-describe}}

El paquete Python y la aplicación web como herramienta abierta de
análisis de sondeos: series continuas de bioturbación a partir de
imágenes, junto con registros multiproxy, intervalos descritos,
clasificación automática de intervalos, mapas por edad y figuras,
organizado en proyectos.

\hypertarget{muxe9todos-que-usa}{%
\subsubsection{Métodos que usa}\label{muxe9todos-que-usa}}

\begin{longtable}[]{@{}
  >{\raggedright\arraybackslash}p{(\columnwidth - 2\tabcolsep) * \real{0.7273}}
  >{\raggedright\arraybackslash}p{(\columnwidth - 2\tabcolsep) * \real{0.2727}}@{}}
\toprule\noalign{}
\begin{minipage}[b]{\linewidth}\raggedright
Método
\end{minipage} & \begin{minipage}[b]{\linewidth}\raggedright
§
\end{minipage} \\
\midrule\noalign{}
\endhead
\bottomrule\noalign{}
\endlastfoot
Escala de profundidad y zona de análisis & §2 \\
Porcentaje por intervalos, multiclase y exclusiones & §3, §8 \\
BI e ii (tramos de la V2, configurables) & §4 \\
Modelo de edad (interpolación lineal) & §5 \\
Realce, detección y edición (aplicación web) & §6--§8 \\
Identificación asistida k-NN (aplicación web) & §10 \\
Formatos de exportación & §12 \\
Resolución de trabajo, proyectos & §16--§17 \\
Registro compuesto y registros externos & §18--§19 \\
Herramienta multiproxy: familias, clasificadores (RF, MLP, k-NN, NB),
validación & §21--§21.5 \\
Mapas: proyecciones, interpolación (IDW, kriging, TIN, RBF\ldots),
exportación & §21.6--§21.6 quater \\
Gráficos de datos (tipo Grapher) y hoja de cálculo & §22.3 bis \\
Organización por proyectos y sondeos, análisis entre proyectos, figuras
& §22 \\
Proyecto de entrenamiento y recorrido guiado & §22.7 \\
\end{longtable}

\hypertarget{quuxe9-se-puede-afirmar-hoy}{%
\subsubsection{Qué se puede afirmar
hoy}\label{quuxe9-se-puede-afirmar-hoy}}

\begin{itemize}
\tightlist
\item
  Integra en una sola herramienta las imágenes, la bioturbación y otros
  registros del mismo sondeo, en profundidad o en edad, con exportación
  vectorial (§18). Demostrado con datos públicos del IODP (sondeo de
  prueba U1385I-1H).
\item
  Reproduce exactamente los resultados publicados en 2020 (prueba de
  regresión, §13.1).
\item
  La aplicación web da los mismos BI e ii que la versión publicada en la
  imagen de ejemplo (100 \%), con una diferencia de porcentaje $\leq$ 0,56
  puntos explicada por los pesos de gris (§3.4, §13.2).
\item
  Es determinista y el archivo de proyecto permite reproducir cualquier
  análisis (§15).
\item
  Tiene pruebas automáticas en integración continua (§13.1).
\item
  Integra en el mismo proyecto la interpretación icnológica, registros
  numéricos de cualquier tipo, intervalos descritos jerárquicos y una
  clasificación automática de intervalos (§21). Sobre esta parte solo
  puede afirmarse que funciona como herramienta; su capacidad de
  detectar eventos reales no está validada (ver artículo 4).
\end{itemize}

\hypertarget{quuxe9-falta-antes-de-enviarlo}{%
\subsubsection{Qué falta antes de
enviarlo}\label{quuxe9-falta-antes-de-enviarlo}}

\begin{itemize}
\tightlist
\item[$\square$]
  Publicar en PyPI y crear la release con DOI en Zenodo.
\item[$\square$]
  Documentación de uso (tutorial con un testigo real) y notebook de
  ejemplo.
\item[$\square$]
  Capa de exclusión también en el paquete Python, para igualar las dos
  implementaciones.
\item[$\square$]
  Alojar la aplicación en GitHub Pages.
\item[$\square$]
  Decidir la licencia definitiva.
\item[$\square$]
  Revisar la bibliografía.
\item[$\square$]
  Dividir \texttt{app.js} en módulos (los revisores de JOSS miran la
  calidad del código).
\item[$\square$]
  Decidir el alcance: un único artículo de software con todo, o dejar la
  parte multiproxy para el artículo 4.
\end{itemize}

\hypertarget{preguntas-probables-de-revisores-joss-revisa-tambiuxe9n-el-cuxf3digo}{%
\subsubsection{Preguntas probables de revisores (JOSS revisa también el
código)}\label{preguntas-probables-de-revisores-joss-revisa-tambiuxe9n-el-cuxf3digo}}

\textbf{¿Qué aporta frente a ImageJ/Fiji, que ya cuenta píxeles?} ImageJ
da un área total. IchAnTool da una \textbf{serie continua por
profundidad o edad} con BI, ii y porcentaje por icnotaxón, lista para
correlacionar con datos geoquímicos. Gestiona varias secciones, zonas
excluidas y modelo de edad en un solo flujo. Ese flujo en ImageJ
requiere macros propias.

\textbf{¿Cómo sé que los resultados son correctos?} Pruebas con casos
calculados a mano y prueba de regresión frente a los resultados
publicados (§13.1). Ambas se ejecutan automáticamente en cada cambio.

\textbf{¿Cómo se eligen los parámetros de realce?} Con los ajustes
publicados (Dorador et al., 2014; Miguez-Salas et al., 2019) o con la
propuesta automática por separabilidad de Otsu (§6), que es propia y se
presenta como punto de partida. No hay en la bibliografía un método
publicado de preselección automática (revisión V12).

\textbf{¿Los ajustes «de artículo» reproducen fielmente los protocolos
publicados?} Reproducen la secuencia y los parámetros publicados donde
existen (§6). Las diferencias están documentadas: los niveles se
calculan por percentiles porque el artículo los elige en cada imagen; la
gamma es automática dentro del rango publicado; el contraste y el brillo
de Photoshop solo tienen equivalentes aproximados; y el CLAHE es por
mosaicos (Zuiderveld, 1994) frente a la ventana deslizante de Fiji. Los
parámetros efectivos de cada sección se guardan en el proyecto.

\textbf{¿Por qué dos implementaciones (Python y JavaScript)?} Python
para análisis por lotes y reproducibilidad en scripts. Web para el
trabajo interactivo sin instalar nada. Comparten las mismas definiciones
y las pruebas comprueban los mismos casos en las dos (§1, §13.1).

\textbf{¿Por qué BI e ii salen solo del porcentaje?} Es la componente
cuantificable de esas escalas. Se documenta que no sustituye la
descripción completa del icnofábrica, y la aplicación guarda el índice
manual del usuario (§4.3).

\begin{center}\rule{0.5\linewidth}{0.5pt}\end{center}

\hypertarget{artuxedculo-2-conjunto-de-datos-anotado-y-segmentaciuxf3n-automuxe1tica}{%
\subsection{Artículo 2 --- Conjunto de datos anotado y segmentación
automática}\label{artuxedculo-2-conjunto-de-datos-anotado-y-segmentaciuxf3n-automuxe1tica}}

\hypertarget{quuxe9-describiruxe1}{%
\subsubsection{Qué describirá}\label{quuxe9-describiruxe1}}

Un conjunto abierto de imágenes de testigos con máscaras de bioturbación
(y de icnotaxones), anotado con la aplicación web, y la evaluación de
métodos automáticos de detección e identificación.

\hypertarget{muxe9todos-que-usa-1}{%
\subsubsection{Métodos que usa}\label{muxe9todos-que-usa-1}}

\begin{longtable}[]{@{}ll@{}}
\toprule\noalign{}
Método & § \\
\midrule\noalign{}
\endhead
\bottomrule\noalign{}
\endlastfoot
Herramientas de anotación (varita, pincel, lazo, exclusión) & §8 \\
Detección clásica: umbral adaptativo + morfología (método base) & §7 \\
Descriptores de objeto & §9 \\
k-NN con rechazo de novedad (método base de identificación) & §10 \\
Segmentación aprendida tipo U-Net (\textbf{por implementar}) & --- \\
Formato de máscara y paleta & §12 \\
\end{longtable}

\hypertarget{quuxe9-se-puede-afirmar-hoy-1}{%
\subsubsection{Qué se puede afirmar
hoy}\label{quuxe9-se-puede-afirmar-hoy-1}}

Nada sobre la precisión en datos reales: todavía no hay datos anotados.
Solo que las herramientas existen y producen máscaras en un formato
abierto.

\begin{itemize}
\tightlist
\item
  La identificación asistida de icnotaxones ya tiene evaluación por
  imagen (V22); falta medirla en testigos independientes anotados por
  expertos.
\end{itemize}

\hypertarget{experimentos-necesarios}{%
\subsubsection{Experimentos necesarios}\label{experimentos-necesarios}}

\begin{enumerate}
\def\labelenumi{\arabic{enumi}.}
\tightlist
\item
  \textbf{Datos:} testigos con imágenes de licencia clara (IODP/ODP son
  públicas; las imágenes tratadas del grupo de la UGR, con permiso).
  Anotar la resolución y el equipo de cada imagen.
\item
  \textbf{Protocolo de anotación escrito}, con criterios de qué es
  traza, qué es artefacto y qué se excluye.
\item
  \textbf{Concordancia entre anotadores:} al menos dos expertos anotan
  el mismo subconjunto. Se mide IoU (intersección sobre unión) y la
  diferencia de porcentaje por intervalo.
\item
  \textbf{División por testigo}, no por recortes: los testigos de prueba
  no se usan para entrenar ni para ajustar parámetros. Evita resultados
  inflados por autocorrelación espacial.
\item
  \textbf{Métricas:}

  \begin{itemize}
  \tightlist
  \item
    IoU y F1 por píxel.
  \item
    Error absoluto del porcentaje por intervalo.
  \item
    Coincidencia de BI e ii.
  \item
    Para la identificación: matriz de confusión y F1 por clase.
  \end{itemize}
\item
  \textbf{Comparaciones:}

  \begin{itemize}
  \tightlist
  \item
    umbral adaptativo (método base);
  \item
    flujo Fiji/ICY de Miguez-Salas et al.~(2019): CLAHE + umbral o
    K-medias + rellenar huecos, ya implementado;
  \item
    tratamiento de Dorador et al.~(2014), ya implementado;
  \item
    U-Net, comparable con Kikuchi y Naruse (2024, Dice 0,892 para
    \emph{Phycosiphon});
  \item
    concordancia entre humanos, que es el techo razonable.
  \end{itemize}
\end{enumerate}

\hypertarget{preguntas-probables-de-revisores}{%
\subsubsection{Preguntas probables de
revisores}\label{preguntas-probables-de-revisores}}

\textbf{¿Por qué un umbral adaptativo como base y no directamente deep
learning?} Porque es reproducible, no necesita datos de entrenamiento y
marca el nivel mínimo que un método aprendido debe superar. Se compara
con U-Net en el mismo conjunto de prueba.

\textbf{¿El acierto del k-NN no está inflado?} El acierto dejando uno
fuera que muestra la aplicación sí puede estarlo (§10, nota al revisor).
En el artículo solo se dará el acierto en \textbf{testigos
independientes}.

\textbf{¿Qué pasa con la dependencia del color y de la cámara?} Se
documentan el equipo y la iluminación de cada imagen y se evalúa el
rendimiento entre fuentes distintas (entrenar con unas, probar con
otras). Los descriptores de tamaño van en unidades físicas (§9).

\textbf{¿Hay trabajos previos con deep learning en bioturbación?} Sí.
Hay que revisarlos y citarlos: el artículo de \emph{Applied Sciences}
11(16):7736 (2021) sobre deep learning en análisis icnológico y la
patente US12142024. La aportación propia es el \textbf{conjunto de datos
abierto}, la herramienta de anotación libre y la comparación con
concordancia humana.

\begin{center}\rule{0.5\linewidth}{0.5pt}\end{center}

\hypertarget{artuxedculo-3-validaciuxf3n-y-aplicaciuxf3n-paleoambiental}{%
\subsection{Artículo 3 --- Validación y aplicación
paleoambiental}\label{artuxedculo-3-validaciuxf3n-y-aplicaciuxf3n-paleoambiental}}

\hypertarget{quuxe9-describiruxe1-1}{%
\subsubsection{Qué describirá}\label{quuxe9-describiruxe1-1}}

\begin{enumerate}
\def\labelenumi{\arabic{enumi}.}
\tightlist
\item
  Cuánto difieren los índices automáticos de la estimación visual de
  varios icnólogos.
\item
  Registros continuos de bioturbación y de icnotaxones en testigos
  largos, correlacionados con indicadores geoquímicos y de temperatura
  (amplía el ejemplo de ODP 977 de 2020).
\end{enumerate}

\hypertarget{muxe9todos-que-usa-2}{%
\subsubsection{Métodos que usa}\label{muxe9todos-que-usa-2}}

\begin{longtable}[]{@{}ll@{}}
\toprule\noalign{}
Método & § \\
\midrule\noalign{}
\endhead
\bottomrule\noalign{}
\endlastfoot
Porcentaje, BI e ii automáticos & §3--§4 \\
BI e ii manuales en la aplicación & §4.3 \\
Porcentaje por icnotaxón e identificación asistida & §3, §10--§11 \\
Modelo de edad & §5 \\
Registro compuesto, registros externos y CSV interpolado & §18--§19 \\
\end{longtable}

\hypertarget{experimentos-necesarios-1}{%
\subsubsection{Experimentos
necesarios}\label{experimentos-necesarios-1}}

\begin{enumerate}
\def\labelenumi{\arabic{enumi}.}
\setcounter{enumi}{-1}
\tightlist
\item
  \textbf{Modelo de edad publicado:} el del sondeo de prueba es
  ilustrativo y no vale para interpretar.
\item
  \textbf{Experimento a ciegas:} de 3 a 5 icnólogos asignan BI e ii
  visuales a los mismos intervalos, sin ver el resultado automático. La
  aplicación ya permite registrarlo (columnas \texttt{BI\_manual} e
  \texttt{ii\_manual}).
\item
  \textbf{Estadística:}

  \begin{itemize}
  \tightlist
  \item
    Concordancia entre observadores (kappa de Cohen ponderado o
    coeficiente de correlación intraclase).
  \item
    Sesgo de cada observador frente al valor automático.
  \item
    Diagrama de Bland-Altman para el porcentaje.
  \end{itemize}
\item
  \textbf{Aplicación:}

  \begin{itemize}
  \tightlist
  \item
    Elegir testigos con buenas imágenes y datos de indicadores ya
    publicados (material de la tesis).
  \item
    Correlación con desfase y prueba de significación que tenga en
    cuenta la autocorrelación de las series.
  \end{itemize}
\end{enumerate}

\hypertarget{preguntas-probables-de-revisores-1}{%
\subsubsection{Preguntas probables de
revisores}\label{preguntas-probables-de-revisores-1}}

\textbf{Si los observadores no coinciden entre sí, ¿con qué se valida el
método?} No se trata de que la máquina «acierte» a un experto. Se mide
cuánta variabilidad hay entre expertos y si el método automático cae
dentro de ella, con la ventaja de ser reproducible.

\textbf{¿Una medida 2D de la cara del testigo representa la
bioturbación?} Es la misma limitación de las estimaciones visuales, que
también se hacen sobre la cara expuesta. Se indica como limitación (§14)
y se mantiene la comparación en las mismas condiciones.

\textbf{¿Cómo afecta la incertidumbre del modelo de edad a las
correlaciones?} Se discute aparte. Si hace falta, se repite el análisis
con modelos de edad alternativos. El programa no propaga esa
incertidumbre (§5).

\begin{center}\rule{0.5\linewidth}{0.5pt}\end{center}

\hypertarget{artuxedculo-4-clasificaciuxf3n-automuxe1tica-de-intervalos-con-registros-multiproxy-descripciones-e-icnologuxeda}{%
\subsection{Artículo 4 --- Clasificación automática de intervalos con
registros multiproxy, descripciones e
icnología}\label{artuxedculo-4-clasificaciuxf3n-automuxe1tica-de-intervalos-con-registros-multiproxy-descripciones-e-icnologuxeda}}

\hypertarget{quuxe9-describiruxe1-2}{%
\subsubsection{Qué describirá}\label{quuxe9-describiruxe1-2}}

Un método reproducible para etiquetar intervalos (eventos como baja
oxigenación, aporte detrítico o corrientes de fondo, pero también
litologías o facies) en muchos sondeos a partir de registros numéricos
(geoquímica, foraminíferos, TOC, sedimentología\ldots) y de intervalos
descritos (litología, turbiditas\ldots), combinando el criterio experto
(intervalos y reglas) con aprendizaje supervisado, y su cartografía por
ventanas de edad. Incluye la icnología (bioturbación, BI, ii,
icnotaxones) como variables.

\hypertarget{muxe9todos-que-usa-3}{%
\subsubsection{Métodos que usa}\label{muxe9todos-que-usa-3}}

\begin{longtable}[]{@{}
  >{\raggedright\arraybackslash}p{(\columnwidth - 2\tabcolsep) * \real{0.7273}}
  >{\raggedright\arraybackslash}p{(\columnwidth - 2\tabcolsep) * \real{0.2727}}@{}}
\toprule\noalign{}
\begin{minipage}[b]{\linewidth}\raggedright
Método
\end{minipage} & \begin{minipage}[b]{\linewidth}\raggedright
§
\end{minipage} \\
\midrule\noalign{}
\endhead
\bottomrule\noalign{}
\endlastfoot
Importación, remuestreo, razones, log y escalado & §21.1--§21.2 \\
Capas de intervalos jerárquicas como variables (0/1) & §21.2--§21.3 \\
Etiquetas por intervalos y reglas & §21.3 \\
Clasificadores definidos por el usuario (familias objetivo, entradas,
modelo) & §21.4 ter \\
Bosque aleatorio, red neuronal (MLP), k-NN, Bayes ingenuo, k-medias &
§21.4, §21.4 bis \\
Clasificación manual en columnas y copia desde columnas automáticas &
§21.4 ter \\
Validación dejando un sondeo fuera & §21.5 \\
Mapas por edad e IDW & §21.6 \\
Bioturbación e índices como variables & §3--§4, §18, §21.1 \\
\end{longtable}

\hypertarget{quuxe9-se-puede-afirmar-hoy-2}{%
\subsubsection{Qué se puede afirmar
hoy}\label{quuxe9-se-puede-afirmar-hoy-2}}

\begin{itemize}
\tightlist
\item
  La implementación funciona y es reproducible (semilla) con datos
  sintéticos de verdad conocida (§21.8).
\item
  Primera prueba con datos reales (§21.8 bis): con registros públicos de
  ODP 977A, los tipos de ORL de la tesis se reconocen mal (acierto
  equilibrado 36 \% por bloques). No se puede afirmar capacidad
  predictiva; faltan los datos geoquímicos propios (XRF) y más sondeos
  etiquetados.
\item
  Con los datos propios (V12, §21.8 ter): acierto equilibrado 17,7 \%
  dejando un sondeo fuera con 7 clases. La transferencia entre sondeos
  con variables y escáneres distintos es el problema principal que un
  artículo tendría que resolver (armonizar variables, más sondeos por
  tipo de variable).
\item
  V26 (misma tesis): el clasificador creado desde los ajustes guardados
  da 18 \% de acierto equilibrado, coherente con la cifra anterior. En
  el proyecto sintético de entrenamiento, bosque aleatorio $\approx$ 91 \% y red
  neuronal $\approx$ 80 \% (verdad conocida; solo demuestra que la cadena
  funciona).
\end{itemize}

\hypertarget{experimentos-necesarios-2}{%
\subsubsection{Experimentos
necesarios}\label{experimentos-necesarios-2}}

\begin{enumerate}
\def\labelenumi{\arabic{enumi}.}
\tightlist
\item
  \textbf{Datos reales con interpretación independiente:} sondeos con
  eventos publicados (p.~ej. sapropeles datados, contornitas, niveles
  detríticos) que no hayan servido para definir las reglas.
\item
  \textbf{Etiquetas solo por intervalos} para entrenar y evaluar. Las
  reglas se evalúan aparte, como «modelo experto», frente a la misma
  verdad.
\item
  \textbf{Validación dejando sondeos enteros fuera} y, si hay
  suficientes, por regiones (bloques espaciales), para medir la
  transferencia a sondeos lejanos.
\item
  \textbf{Métricas:} acierto equilibrado, F1 por clase y matriz de
  confusión; incertidumbre por remuestreo de sondeos (bootstrap por
  sondeo).
\item
  \textbf{Comparaciones:} reglas expertas solas; bosque aleatorio; red
  neuronal; k-NN; Bayes; y el efecto de añadir la icnología (con y sin
  bioturbación/índices) y las capas descritas (con y sin litología) como
  variables. Vigilar que una capa descrita no sea casi la misma
  información que el objetivo (fuga de información), y declarar qué
  capas se usaron.
\item
  \textbf{Importancia por permutación} (por implementar) además de MDI.
\item
  \textbf{Sensibilidad} al paso de remuestreo, al escalado (global
  frente a por sondeo) y a la probabilidad mínima.
\item
  \textbf{Modelos de edad:} repetir los mapas con modelos alternativos o
  con la incertidumbre publicada de cada sondeo.
\item
  \textbf{Prueba final independiente:} reservar 2--3 sondeos que no se
  usan para elegir modelo ni ajustes, y dar en el artículo el acierto en
  ellos.
\end{enumerate}

\hypertarget{preguntas-probables-de-revisores-2}{%
\subsubsection{Preguntas probables de
revisores}\label{preguntas-probables-de-revisores-2}}

\textbf{¿El acierto no está inflado por la autocorrelación de las
series?} Por eso la validación deja sondeos enteros fuera (o bloques
contiguos si solo hay uno) y nunca mezcla muestras vecinas entre
entrenamiento y prueba (§21.5).

\textbf{Si las etiquetas salen de reglas, ¿el modelo no aprende solo las
reglas?} Sí, y así se declara (§21.8). Para la evaluación se usan
intervalos con interpretación independiente; las reglas solo sirven para
proponer etiquetas o como referencia experta.

\textbf{¿Por qué no deep learning?} El programa incluye una red neuronal
(perceptrón multicapa con parada temprana, §21.4 bis) y se compara con
el bosque aleatorio, k-NN y Bayes en el mismo esquema de validación. Con
pocas decenas o cientos de intervalos etiquetados, los árboles suelen
ser más robustos y dan una medida de importancia; redes más profundas
(recurrentes, convolucionales sobre la serie) necesitarían muchos más
sondeos etiquetados de los disponibles.

\textbf{Probasteis varios modelos: ¿no estáis eligiendo el que mejor
sale por azar?} Se informa de todos los modelos probados con la misma
validación y el acierto final se da en sondeos reservados que no
intervinieron en la elección (experimento 9).

\textbf{¿La interpolación del mapa es fiable?} La interpolación es
visual y exploratoria (§21.6, §21.6 ter): ningún método estima la
incertidumbre de forma útil con pocos sondeos (el kriging da una
varianza, pero su variograma con \textasciitilde10 sondeos es muy
incierto) y ninguno tiene en cuenta barreras geográficas. Las
conclusiones se basan en los valores de cada sondeo, que también se
exportan. Los fondos de los mapas (Esri, GEBCO; §21.6 bis) son solo
contexto visual: no entran en ningún cálculo, y su fuente debe citarse
en el pie de la figura (el programa la escribe en el mapa).

\textbf{¿Cómo entran datos categóricos como la litología en un modelo
numérico?} Como columnas 0/1 por categoría, sin escalar (§21.2). La
ausencia de descripción vale 0, igual que «no presente»; se declara como
limitación y conviene usar solo sondeos con la capa descrita entera.

\textbf{¿Cómo se combinan proxies medidos con técnicas o calibraciones
distintas?} Con el z-score por sondeo (§21.2) se eliminan diferencias de
nivel; se documenta qué escalado se usó y se muestra su efecto
(experimento 7).

\begin{center}\rule{0.5\linewidth}{0.5pt}\end{center}

\hypertarget{artuxedculo-5-nota-breve-en-espauxf1ol-geogaceta}{%
\subsection{Artículo 5 --- Nota breve en español
(Geogaceta)}\label{artuxedculo-5-nota-breve-en-espauxf1ol-geogaceta}}

\hypertarget{quuxe9-describiruxe1-3}{%
\subsubsection{Qué describirá}\label{quuxe9-describiruxe1-3}}

Presentación de CorAnTool a la comunidad geológica española: qué hace,
para quién es y un ejemplo con un testigo de datos públicos (IODP
U1385I-1H, margen ibérico) y el proyecto de entrenamiento. Sin
resultados inéditos.

\hypertarget{quuxe9-se-puede-afirmar-hoy-3}{%
\subsubsection{Qué se puede afirmar
hoy}\label{quuxe9-se-puede-afirmar-hoy-3}}

Lo mismo que en el artículo 1 (funcionalidad y reproducibilidad), sin
cifras de validación nuevas.

\hypertarget{quuxe9-falta}{%
\subsubsection{Qué falta}\label{quuxe9-falta}}

\begin{itemize}
\tightlist
\item[$\square$]
  Release con DOI y aplicación accesible en línea (GitHub Pages).
\item[$\square$]
  Figura: imagen tratada + columna compuesta + mapa de ejemplo.
\item[$\square$]
  Comprobar formato y extensión en las normas de la revista.
\end{itemize}

\hypertarget{preguntas-probables-de-revisores-3}{%
\subsubsection{Preguntas probables de
revisores}\label{preguntas-probables-de-revisores-3}}

\textbf{¿Qué aporta frente a hacerlo a mano en Photoshop/ImageJ y una
hoja de cálculo?} Un flujo único, reproducible y gratuito desde la
imagen hasta series, mapas y figuras, con el proyecto guardado para
repetir el análisis.

\begin{center}\rule{0.5\linewidth}{0.5pt}\end{center}

\hypertarget{lista-de-comprobaciuxf3n-antes-de-enviar-cualquier-artuxedculo}{%
\subsection{Lista de comprobación antes de enviar cualquier
artículo}\label{lista-de-comprobaciuxf3n-antes-de-enviar-cualquier-artuxedculo}}

\begin{itemize}
\tightlist
\item[$\square$]
  Cada método citado en el texto tiene su sección en
  \texttt{DOCUMENTACION\_TECNICA.md}, y coincide con el código de la
  versión citada.
\item[$\square$]
  Versión del software citada con DOI de Zenodo.
\item[$\square$]
  Parámetros usados guardados (archivo de proyecto) y publicados como
  material suplementario.
\item[$\square$]
  Datos y máscaras depositados (Zenodo o PANGAEA) con licencia.
\item[$\square$]
  Limitaciones declaradas (§14).
\item[$\square$]
  Referencias comprobadas contra los originales.
\item[$\square$]
  Coautorías y agradecimientos acordados.
\end{itemize}

\newpage

\hypertarget{borrador-1-software-joss-en-ingluxe9s}{%
\section{Borrador 1 --- Software (JOSS, en
inglés)}\label{borrador-1-software-joss-en-ingluxe9s}}

\hypertarget{summary}{%
\subsection{Summary}\label{summary}}

Bioturbation intensity is a key parameter of the ichnofabric approach
and is routinely estimated in sediment cores with semi-quantitative
visual schemes such as the Bioturbation Index (Reineck 1963; Taylor y
Goldring 1993) and the ichnofabric index (Droser y Bottjer 1986). Visual
estimation is subjective, slow, and usually discontinuous along the
core. IchAnTool converts segmented core images into continuous, depth-
or age-referenced series of bioturbated-surface percentage, per-class
coverage (e.g.~per ichnotaxon), and the corresponding semi-quantitative
indices, ready to be correlated with geochemical, mineralogical, or
logging data.

\hypertarget{statement-of-need}{%
\subsection{Statement of need}\label{statement-of-need}}

Digital image treatment has improved the visibility and quantification
of trace fossils in cores (Dorador y Rodríguez-Tovar 2014, 2018), but
turning treated images into continuous data series remained a manual,
repetitive task. A first version (IchAnTool), a single script, was
introduced in Casanova-Arenillas, Rodríguez-Tovar, y Martínez-Ruiz
(2020). CorAnTool is a complete rewrite: an installable, tested Python
package with a command-line interface and API, and a browser application
that needs no installation, aimed at sedimentologists, ichnologists,
palaeoceanographers, and core-analysis practitioners in academia and
industry. {[}TODO: one paragraph on existing alternatives and why
IchAnTool is needed.{]}

\hypertarget{functionality}{%
\subsection{Functionality}\label{functionality}}

\begin{itemize}
\tightlist
\item
  Input: treated images (PNG, TIFF, JPEG) with bioturbation painted in a
  single colour, colour-coded classes, or integer label masks.
\item
  Row-wise bioturbated percentage, aggregated into user-defined depth
  bins.
\item
  Per-class percentages for colour-coded ichnotaxa.
\item
  Configurable index schemes (BI and ii provided).
\item
  Depth-to-age conversion with user age models.
\item
  CSV/Excel export, quick-look plots, and a modular segmentation
  interface for future learned models.
\item
  Results reproduce those of the published version exactly (regression
  test).
\item
  Browser application: projects with several cores; numeric records of
  any kind (XRF, foraminifera, TOC, colour) in depth or age; families of
  described intervals classified by hand; user-defined classifiers
  (random forest, multilayer perceptron, k-nearest neighbours, naive
  Bayes, k-means) validated leaving whole cores out; maps by age window
  with several projections and interpolation methods; Grapher-like
  plots; publication figures; a synthetic training project with a guided
  tour.
\end{itemize}

\hypertarget{implementation}{%
\subsection{Implementation}\label{implementation}}

Both implementations share the same definitions. A label mask (0 =
matrix, 1..n = objects or classes) is restricted to an analysis area.
The bioturbated percentage of each depth interval of \emph{N} pixel rows
is the ratio of bioturbated to valid pixels. Excluded pixels (cracks,
voids, samples; currently in the web application only) leave both terms.
Each interval is assigned the depth of its centre, interpolated linearly
between the top and bottom depths of the analysis area. Percentages map
to BI and ii through half-open tiers, and ages come from linear
interpolation of a user age model.

The web application adds four steps before the mask is final:

\begin{itemize}
\tightlist
\item
  image enhancement: levels, gamma, contrast, saturation, histogram
  equalisation and unsharp-mask local contrast;
\item
  candidate detection: adaptive mean-C thresholding, 3 × 3 morphological
  opening/closing and 8-connected components with a minimum size;
\item
  interactive editing: colour region growing, brush, lasso and an
  exclusion layer;
\item
  assisted ichnotaxon identification: a distance-weighted
  \emph{k}-nearest-neighbour classifier (\emph{k} = 5) on twelve
  standardised shape, size (physical units) and colour descriptors, with
  a confidence threshold, a novelty gate based on the leave-one-out
  nearest-neighbour distance and exportable models.
\end{itemize}

All computations are deterministic. Project files store the original
image and every parameter. A full technical reference, with formulas and
default values, is distributed with the code (\texttt{docs/}).

{[}TODO: figure --- example core, mask and output curves
(\texttt{examples/}).{]}

\hypertarget{acknowledgements}{%
\subsection{Acknowledgements}\label{acknowledgements}}

{[}TODO{]}

\hypertarget{references}{%
\subsection{References}\label{references}}

\newpage

\hypertarget{borrador-2-dataset-y-segmentaciuxf3n}{%
\section{Borrador 2 --- Dataset y
segmentación}\label{borrador-2-dataset-y-segmentaciuxf3n}}

\hypertarget{paper-2-open-annotated-dataset-and-automatic-segmentation-of-bioturbation-in-core-images}{%
\subsection{Paper 2 --- Open annotated dataset and automatic
segmentation of bioturbation in core
images}\label{paper-2-open-annotated-dataset-and-automatic-segmentation-of-bioturbation-in-core-images}}

\textbf{Target:} Earth System Science Data (data focus) or Computers \&
Geosciences / AI in Geosciences (method focus). Decide once the dataset
size and model results are known.

\hypertarget{working-title}{%
\subsubsection{Working title}\label{working-title}}

An open, annotated dataset of sediment-core images for automatic
segmentation of bioturbation and ichnotaxa {[}TODO{]}

\hypertarget{provisional-abstract}{%
\subsubsection{Provisional abstract}\label{provisional-abstract}}

Quantifying bioturbation in cores still depends on manual selection of
trace fossils on images, the main bottleneck of image-based ichnofabric
analysis. We present {[}N{]} core images from {[}sites/expeditions{]}
with expert masks of bioturbation {[}and ichnotaxa{]}, and train
{[}model{]} to segment them automatically. Automatic masks are compared
with expert masks and with visual BI/ii estimates. Data, code
(IchAnTool) and trained models are openly released. {[}TODO: results{]}

\hypertarget{outline}{%
\subsubsection{Outline}\label{outline}}

\begin{enumerate}
\def\labelenumi{\arabic{enumi}.}
\tightlist
\item
  Introduction --- subjectivity and cost of manual selection; image
  treatment state of the art.
\item
  Related work --- recent deep-learning approaches to bioturbation
  {[}TODO: review, incl. Applied Sciences 11(16):7736 (2021) and patent
  US12142024; verify authors{]}.
\item
  Data --- sources (ODP/IODP, UGR collections), resolution, licences,
  annotation protocol, inter-annotator agreement.
\item
  Methods --- annotation tool, model, training/validation split by core
  (not by tile).
\item
  Results --- IoU/F1, percentage error per bin, BI/ii agreement.
\item
  Usage --- integration with IchAnTool (\texttt{segmentation} module).
\item
  Data availability --- Zenodo/PANGAEA DOI.
\end{enumerate}

\hypertarget{needs}{%
\subsubsection{Needs}\label{needs}}

\begin{itemize}
\tightlist
\item[$\square$]
  Permission and licence for the images (IODP images are public; UGR
  treated images: ask).
\item[$\square$]
  Annotation by $\geq$2 experts on a subset to measure agreement.
\item[$\square$]
  GPU for training (a single consumer GPU is enough for a U-Net).
\end{itemize}

\newpage

\hypertarget{borrador-3-validaciuxf3n-y-aplicaciuxf3n}{%
\section{Borrador 3 --- Validación y
aplicación}\label{borrador-3-validaciuxf3n-y-aplicaciuxf3n}}

\hypertarget{paper-3-validation-and-palaeoenvironmental-application}{%
\subsection{Paper 3 --- Validation and palaeoenvironmental
application}\label{paper-3-validation-and-palaeoenvironmental-application}}

\textbf{Target:} Palaeogeography, Palaeoclimatology, Palaeoecology.
Alternatives: Ichnos, Marine Geology, Sedimentary Geology.

\hypertarget{working-title-1}{%
\subsubsection{Working title}\label{working-title-1}}

Continuous, objective bioturbation records from core images: validation
against visual indices and application to {[}region/period{]} {[}TODO{]}

\hypertarget{idea}{%
\subsubsection{Idea}\label{idea}}

\begin{itemize}
\tightlist
\item
  Validation: compare IchAnTool percentages with blind visual BI/ii
  estimates by several ichnologists on the same intervals; quantify
  subjectivity and bias of visual schemes.
\item
  Application: continuous bioturbation and per-ichnotaxon records along
  long cores, correlated with geochemical/redox and SST proxies
  (e.g.~Western Mediterranean during the Last Glacial Cycle; extends the
  ODP 977 example of Casanova-Arenillas et al., 2020).
\end{itemize}

\hypertarget{outline-1}{%
\subsubsection{Outline}\label{outline-1}}

\begin{enumerate}
\def\labelenumi{\arabic{enumi}.}
\tightlist
\item
  Introduction
\item
  Material --- cores, images, age models, proxy data
\item
  Methods --- IchAnTool workflow, observer experiment, statistics
\item
  Results --- agreement and bias; continuous records
\item
  Discussion --- what continuous records reveal that discrete estimates
  miss
\item
  Conclusions
\end{enumerate}

\hypertarget{needs-1}{%
\subsubsection{Needs}\label{needs-1}}

\begin{itemize}
\tightlist
\item[$\square$]
  Choose cores with good images and existing proxy data (thesis
  material).
\item[$\square$]
  Recruit 3--5 observers for the blind test.
\end{itemize}

\newpage

\hypertarget{borrador-4-clasificaciuxf3n-multiproxy}{%
\section{Borrador 4 --- Clasificación
multiproxy}\label{borrador-4-clasificaciuxf3n-multiproxy}}

\hypertarget{paper-4-automatic-classification-of-core-intervals-from-multiproxy-records}{%
\subsection{Paper 4 --- Automatic classification of core intervals from
multiproxy
records}\label{paper-4-automatic-classification-of-core-intervals-from-multiproxy-records}}

\textbf{Target:} Computers \& Geosciences. Alternatives: Artificial
Intelligence in Geosciences, Applied Computing and Geosciences; Marine
Geology if the emphasis is the palaeoceanographic result.

\hypertarget{working-title-2}{%
\subsubsection{Working title}\label{working-title-2}}

Reproducible classification of core intervals (sapropels, organic-rich
layers, turbidites\ldots) from multiproxy records with user-defined
machine-learning classifiers {[}TODO{]}

\hypertarget{provisional-abstract-1}{%
\subsubsection{Provisional abstract}\label{provisional-abstract-1}}

Interval interpretation in cores (e.g.~low-oxygen events, detrital
pulses, climate stages) is done by experts core by core and is hard to
reproduce across many cores and laboratories. We present a workflow in
CorAnTool in which experts define families of intervals and their
classes, label some cores, and train classifiers (random forest,
multilayer perceptron, k-nearest neighbours, naive Bayes) on resampled
multiproxy records (XRF, foraminifera, TOC, colour, bioturbation).
Models are validated leaving whole cores out and tested on {[}N{]} cores
withheld from model selection. {[}TODO: results --- balanced accuracy
per family, per class F1, transfer between basins, effect of adding
ichnology.{]}

\hypertarget{outline-2}{%
\subsubsection{Outline}\label{outline-2}}

\begin{enumerate}
\def\labelenumi{\arabic{enumi}.}
\tightlist
\item
  Introduction --- expert interval interpretation; reproducibility
  across cores; prior ML in core logging {[}TODO review{]}.
\item
  Data --- cores, proxies, age models, described intervals (independent
  of the rules), harmonisation of variables.
\item
  Methods --- resampling and scaling (§21.2); families and labels
  (§21.3); classifiers (§21.4--§21.4 ter); validation leaving one core
  out (§21.5); held-out cores for the final test; permutation importance
  {[}TODO implement{]}.
\item
  Results --- comparison of models and inputs; maps by age window
  (§21.6) as illustration only.
\item
  Discussion --- transfer between cores with different scanners; leakage
  between families; limits.
\item
  Software and data availability --- CorAnTool release (DOI), project
  files as supplementary material.
\end{enumerate}

\hypertarget{needs-2}{%
\subsubsection{Needs}\label{needs-2}}

\begin{itemize}
\tightlist
\item[$\square$]
  Cores with independently described events (published sapropels/ORL,
  turbidites) not used to define rules.
\item[$\square$]
  Harmonised variable set across cores (same XRF ratios or calibrated
  values).
\item[$\square$]
  Held-out cores option in the interface (to avoid model-selection
  bias).
\item[$\square$]
  Permutation importance (also valid for the neural network).
\end{itemize}

\newpage

\hypertarget{borrador-5-nota-en-geogaceta}{%
\section{Borrador 5 --- Nota en
Geogaceta}\label{borrador-5-nota-en-geogaceta}}

\hypertarget{artuxedculo-5-nota-breve-corantool-una-herramienta-abierta-para-el-anuxe1lisis-de-sondeos}{%
\subsection{Artículo 5 --- Nota breve: CorAnTool, una herramienta
abierta para el análisis de
sondeos}\label{artuxedculo-5-nota-breve-corantool-una-herramienta-abierta-para-el-anuxe1lisis-de-sondeos}}

\textbf{Revista:} Geogaceta (Sociedad Geológica de España). Alternativa:
Spanish Journal of Palaeontology.

\hypertarget{tuxedtulo-provisional}{%
\subsubsection{Título provisional}\label{tuxedtulo-provisional}}

CorAnTool: una herramienta abierta para cuantificar la bioturbación e
integrar datos multiproxy en sondeos {[}TODO{]}

\hypertarget{idea-1}{%
\subsubsection{Idea}\label{idea-1}}

Presentar el programa a la comunidad española con un ejemplo
reproducible sobre datos públicos (IODP U1385I-1H, margen ibérico) y el
proyecto de entrenamiento: imagen tratada $\rightarrow$ porcentaje de bioturbación,
BI e ii $\rightarrow$ columna compuesta con registros geoquímicos $\rightarrow$ familias de
intervalos y clasificación $\rightarrow$ mapa por edad.

\hypertarget{esquema}{%
\subsubsection{Esquema}\label{esquema}}

\begin{enumerate}
\def\labelenumi{\arabic{enumi}.}
\tightlist
\item
  Introducción: subjetividad de los índices visuales y falta de
  herramientas abiertas que integren imagen y datos.
\item
  La herramienta: flujo, formatos, reproducibilidad (proyecto guardado,
  semilla), aplicación web sin instalación.
\item
  Ejemplo con datos públicos (figura única con tres paneles).
\item
  Disponibilidad: código, DOI, licencia, aplicación en línea.
\end{enumerate}

\hypertarget{pendiente}{%
\subsubsection{Pendiente}\label{pendiente}}

\begin{itemize}
\tightlist
\item[$\square$]
  Normas de extensión y formato de la revista.
\item[$\square$]
  Figura y aplicación publicada en GitHub Pages.
\end{itemize}

\newpage

\hypertarget{revistas-candidatas}{%
\section{Revistas candidatas}\label{revistas-candidatas}}

No se incluyen factores de impacto ni APC: cambian cada año. Verificar
en la web de cada revista antes de enviar.

\hypertarget{software-y-muxe9todos-computacionales}{%
\subsection{Software y métodos
computacionales}\label{software-y-muxe9todos-computacionales}}

\begin{longtable}[]{@{}
  >{\raggedright\arraybackslash}p{(\columnwidth - 6\tabcolsep) * \real{0.2500}}
  >{\raggedright\arraybackslash}p{(\columnwidth - 6\tabcolsep) * \real{0.3056}}
  >{\raggedright\arraybackslash}p{(\columnwidth - 6\tabcolsep) * \real{0.2222}}
  >{\raggedright\arraybackslash}p{(\columnwidth - 6\tabcolsep) * \real{0.2222}}@{}}
\toprule\noalign{}
\begin{minipage}[b]{\linewidth}\raggedright
Revista
\end{minipage} & \begin{minipage}[b]{\linewidth}\raggedright
Editorial
\end{minipage} & \begin{minipage}[b]{\linewidth}\raggedright
Acceso
\end{minipage} & \begin{minipage}[b]{\linewidth}\raggedright
Encaje
\end{minipage} \\
\midrule\noalign{}
\endhead
\bottomrule\noalign{}
\endlastfoot
Journal of Open Source Software (JOSS) & Open Journals & Abierto, sin
coste & Artículo corto + revisión del código en GitHub. Ideal para
IchAnTool 3. \\
Computers \& Geosciences & Elsevier & Híbrido & Donde se publicó la v2.1
(2020). Bueno para un método nuevo (segmentación automática). \\
Applied Computing and Geosciences & Elsevier & Abierto & Revista hermana
de C\&G, orientada a aplicaciones y software. \\
Artificial Intelligence in Geosciences & KeAi & Abierto & Encaje directo
si el aporte principal es el modelo de deep learning. \\
Geoscientific Instrumentation, Methods and Data Systems & Copernicus &
Abierto & Métodos y sistemas de datos; alternativa a C\&G. \\
\end{longtable}

\hypertarget{datos}{%
\subsection{Datos}\label{datos}}

\begin{longtable}[]{@{}
  >{\raggedright\arraybackslash}p{(\columnwidth - 6\tabcolsep) * \real{0.2500}}
  >{\raggedright\arraybackslash}p{(\columnwidth - 6\tabcolsep) * \real{0.3056}}
  >{\raggedright\arraybackslash}p{(\columnwidth - 6\tabcolsep) * \real{0.2222}}
  >{\raggedright\arraybackslash}p{(\columnwidth - 6\tabcolsep) * \real{0.2222}}@{}}
\toprule\noalign{}
\begin{minipage}[b]{\linewidth}\raggedright
Revista
\end{minipage} & \begin{minipage}[b]{\linewidth}\raggedright
Editorial
\end{minipage} & \begin{minipage}[b]{\linewidth}\raggedright
Acceso
\end{minipage} & \begin{minipage}[b]{\linewidth}\raggedright
Encaje
\end{minipage} \\
\midrule\noalign{}
\endhead
\bottomrule\noalign{}
\endlastfoot
Earth System Science Data (ESSD) & Copernicus & Abierto & Dataset
anotado de imágenes de testigos y máscaras de bioturbación. \\
Scientific Data & Springer Nature & Abierto & Data descriptor; alta
visibilidad fuera de geología. \\
Geoscience Data Journal & Wiley & Abierto & Alternativa para datasets
geocientíficos. \\
\end{longtable}

\hypertarget{icnologuxeda-sedimentologuxeda-y-paleoambiente}{%
\subsection{Icnología, sedimentología y
paleoambiente}\label{icnologuxeda-sedimentologuxeda-y-paleoambiente}}

\begin{longtable}[]{@{}
  >{\raggedright\arraybackslash}p{(\columnwidth - 6\tabcolsep) * \real{0.2500}}
  >{\raggedright\arraybackslash}p{(\columnwidth - 6\tabcolsep) * \real{0.3056}}
  >{\raggedright\arraybackslash}p{(\columnwidth - 6\tabcolsep) * \real{0.2222}}
  >{\raggedright\arraybackslash}p{(\columnwidth - 6\tabcolsep) * \real{0.2222}}@{}}
\toprule\noalign{}
\begin{minipage}[b]{\linewidth}\raggedright
Revista
\end{minipage} & \begin{minipage}[b]{\linewidth}\raggedright
Editorial
\end{minipage} & \begin{minipage}[b]{\linewidth}\raggedright
Acceso
\end{minipage} & \begin{minipage}[b]{\linewidth}\raggedright
Encaje
\end{minipage} \\
\midrule\noalign{}
\endhead
\bottomrule\noalign{}
\endlastfoot
Ichnos & Taylor \& Francis & Híbrido & Revista de referencia en
icnología; validación frente a BI/ii visual. \\
Palaeogeography, Palaeoclimatology, Palaeoecology & Elsevier & Híbrido &
Aplicación paleoambiental (p.~ej. Último Ciclo Glacial en el
Mediterráneo). \\
Marine Geology & Elsevier & Híbrido & Aplicación en testigos marinos
(IODP/ODP). \\
Sedimentary Geology & Elsevier & Híbrido & Enfoque
sedimentológico/icnofábrica. \\
Palaios & SEPM & Híbrido & Paleoecología e icnología. \\
Lethaia & Scandinavian University Press & Abierto & Paleontología e
icnología. \\
Marine and Petroleum Geology / AAPG Bulletin & Elsevier / AAPG & Híbrido
& Si se orienta a reservorios y caracterización de testigos en
industria. \\
\end{longtable}

\hypertarget{divulgaciuxf3n-tuxe9cnica-en-espauxf1ol}{%
\subsection{Divulgación técnica en
español}\label{divulgaciuxf3n-tuxe9cnica-en-espauxf1ol}}

\begin{longtable}[]{@{}
  >{\raggedright\arraybackslash}p{(\columnwidth - 4\tabcolsep) * \real{0.3214}}
  >{\raggedright\arraybackslash}p{(\columnwidth - 4\tabcolsep) * \real{0.3929}}
  >{\raggedright\arraybackslash}p{(\columnwidth - 4\tabcolsep) * \real{0.2857}}@{}}
\toprule\noalign{}
\begin{minipage}[b]{\linewidth}\raggedright
Revista
\end{minipage} & \begin{minipage}[b]{\linewidth}\raggedright
Editorial
\end{minipage} & \begin{minipage}[b]{\linewidth}\raggedright
Encaje
\end{minipage} \\
\midrule\noalign{}
\endhead
\bottomrule\noalign{}
\endlastfoot
Geogaceta & Sociedad Geológica de España & Nota corta presentando la
herramienta a la comunidad española. \\
Spanish Journal of Palaeontology & Sociedad Española de Paleontología &
Aplicaciones icnológicas. \\
\end{longtable}

\hypertarget{criterios-para-elegir}{%
\subsection{Criterios para elegir}\label{criterios-para-elegir}}

\begin{enumerate}
\def\labelenumi{\arabic{enumi}.}
\tightlist
\item
  ¿El aporte es el software, los datos, el método o la interpretación
  geológica? Eso decide la familia de revistas.
\item
  Acceso abierto: aumenta adopción de la herramienta; JOSS y Copernicus
  no cobran o cobran poco.
\item
  Plazo: JOSS y las revistas de datos suelen ser más rápidas que las de
  paleoambiente.
\end{enumerate}

\hypertarget{refs}{}
\begin{CSLReferences}{1}{0}
\leavevmode\vadjust pre{\hypertarget{ref-casanova2020}{}}%
Casanova-Arenillas, Santiago, Francisco J. Rodríguez-Tovar, y Francisca
Martínez-Ruiz. 2020. {«Applied ichnology in sedimentary geology: Python
scripts as a method to automatize ichnofabric analysis in marine core
images»}. \emph{Computers \& Geosciences} 136: 104407.
\url{https://doi.org/10.1016/j.cageo.2020.104407}.

\leavevmode\vadjust pre{\hypertarget{ref-dorador2014}{}}%
Dorador, Javier, y Francisco J. Rodríguez-Tovar. 2014. {«Quantitative
estimation of bioturbation based on digital image analysis»}.
\emph{Marine Geology} 349: 55-60.

\leavevmode\vadjust pre{\hypertarget{ref-dorador2018}{}}%
---------. 2018. {«High-resolution image treatment in ichnological core
analysis: Initial steps, advances and prospects»}. \emph{Earth-Science
Reviews} 177: 226-37.

\leavevmode\vadjust pre{\hypertarget{ref-droser1986}{}}%
Droser, Mary L., y David J. Bottjer. 1986. {«A semiquantitative field
classification of ichnofabric»}. \emph{Journal of Sedimentary Petrology}
56: 558-59.

\leavevmode\vadjust pre{\hypertarget{ref-reineck1963}{}}%
Reineck, H.-E. 1963. {«Sedimentgef{ü}ge im Bereich der s{ü}dlichen
Nordsee»}. \emph{Abhandlungen der Senckenbergischen Naturforschenden
Gesellschaft} 505: 1-138.

\leavevmode\vadjust pre{\hypertarget{ref-taylor1993}{}}%
Taylor, A. M., y R. Goldring. 1993. {«Description and analysis of
bioturbation and ichnofabric»}. \emph{Journal of the Geological Society}
150: 141-48.

\end{CSLReferences}

\end{document}
